\chapter*{前言}
\addcontentsline{toc}{chapter}{前言}
\markboth{前言}{前言}

本文集是 Langlands 纲领暑期学校的会议录。暑期学校于 2022 年 7 月 11 日至 29 日在法国高等科学研究所（Institut des Hautes Études Scientifiques，IHES）举办，为期三周。学校由 IHES、美国国家科学基金会（NSF）与 Clay 数学研究所共同资助。报告人共 31 位，参加者逾百人，多为博士生与博士后。

本次暑期学校纪念 1977 年 7 月 11 日至 8 月 5 日在俄勒冈州 Corvallis 举行的 AMS 暑期研究学会「自守形式、表示与 L-函数」45 周年。那次会议录（下称 Corvallis 会议录）在同一 AMS 丛书中出版，此后一直是希望进入这一迷人领域的学生与研究者的入门读物与关键参考文献。此后当然还有许多同类的教学会议与会议录，例如 1988 年 7 月在 Ann Arbor 举行的「自守形式、志村簇与 L-函数」会议（由 Academic Press 出版），1996 年 3 月在爱丁堡举行的「表示论与自守形式」会议（亦由本 AMS 丛书出版），以及 2003 年 6 月在 Fields 研究所举行的 Clay 数学研究所「调和分析、迹公式与志村簇」会议（由 Clay 与 AMS 出版）。然而，正如我们在下文以致敬的方式简要回顾的那样，Corvallis 会议录的广度、深度与影响，在 Langlands 纲领中至今仍未被超越。

Corvallis 暑期研究学会举办于 Langlands 致 Weil 的那封富有远见的信之后约十年，其作用是为学科的进一步发展奠定坚实基础。会议录中的文章给出了基本概念（如自守表示、L-群、自守 L-函数、L-包）的定义，并引入了基本工具（如 Bruhat–Tits 建筑物、Eisenstein 级数、迹公式、志村簇）。Corvallis 会议录把这些思想传播到少数专家之外，使它们如今被视为数论与表示论广大领域中的标准工具。它们还激发了大量后续发展；在此后 45 年间，Corvallis 提出的许多预期已经实现。以下是一份不完全的清单：

\begin{itemize}[leftmargin=2em,itemsep=0.15em]
\item 一般线性群局部 Langlands 猜想（LLC）的证明，以及涵盖所有约化群（包括非拟分裂群）的确定性 LLC 的表述。
\item 描述离散自守谱的 Arthur 猜想及其局部对应物的表述；阿基米德情形及更一般情形下局部 Arthur 包的构造。
\item 逆定理与 Langlands–Shahidi 方法的发展，它们被用来证明 Sym$^3$ 与 Sym$^4$ 函子性提升，以及从典型群到 $\mathrm{GL}_N$ 的整体通用尖点表示的提升。
\item （扭）内窥理论的发展。
\item （扭）迹公式的稳定化（包括基本引理的证明），及其在（拟分裂）典型群自守离散谱分类中的应用。
\item 对志村簇上同调的理解，及其在伽罗瓦表示构造中的应用。
\item 模性提升定理的发展。这些定理支撑了数论中一些最重要的进展，从费马大定理到 Sato–Tate 猜想，再到全纯模形式的对称幂函子性。
\item 在函数域上，一般线性群的 Langlands 猜想以及一般约化群的 Langlands 参数化的证明。
\end{itemize}

这些成就（以及其他成就）见证了 Corvallis 会议录的巨大影响。

鉴于 Langlands 纲领今日的规模，2022 年 IHES 暑期学校若试图效仿 Corvallis、覆盖纲领的全部内容，将是徒劳的；在一个新方向与新联系不断涌现的领域里，试图界定纲领的边界也不合适。我们最初的意图是聚焦于 Langlands 的函子性猜想。但在暑期学校之前的几年里，出现了若干统一的主题与新的联系，它们从根本上改变了我们对 Langlands 纲领的思考方式。因此，暑期学校的组织者（亦即本文集各卷的编者）放弃了最初的设想，邀请了范围广泛的报告人，请他们着重介绍这些新的联系与发展，以帮助克服不同子领域之间时而存在的语言障碍。我们希望这一目标已经达到，也希望贯穿 Langlands 纲领的那些联系线索，在本文集中比以往任何时候都更加清晰可见。

在选择收录内容时，一个主题的相对重要性（按组织者的判断）并非唯一标准。若干重要方向被压缩或略去，或是因为它们已在其他会议与文集中得到充分介绍，或是因为若不对暑期学校的广度作严重牺牲，就无法公正地处理它们。特别地，暑期学校并未着重介绍上文所列过去 45 年间的某些成就。举几个例子：模性方面的进展在会上只得到极简略的覆盖，全纯模形式的对称幂函子性则完全没有涉及；来自逆定理与 Langlands–Shahidi 方法的进展也未真正讨论。几何 Langlands 纲领同样没有单独设讲，但其思想渗透在本文集各卷的许多文章中。我们还可以继续列举未被收录或未获充分覆盖的重要主题。

有鉴于此，尽管本文集记录了 2022 年至本卷出版之间 Langlands 纲领某些方向的最前沿状况，我们仍请读者以宽广而开放的眼光看待这一领域，去探索那些被略去或尚待发现的主题。

下面按卷说明我们在本文集中着重强调的方向。

\begin{itemize}[leftmargin=2em,itemsep=0.15em]
\item 第一卷（第 1、2 部分）简要总结了过去数十年间一直处于、并将继续处于 Langlands 纲领前沿的主题的现状：不可约表示、内窥纲领、志村簇、局部对称空间、shtuka 的模空间及其上同调。
\item 第二卷（第 3 部分）致力于 Langlands 纲领的几何化与范畴化。这里的焦点从单个表示转向表示的族、参数空间上的几何结构，以及这些结构对局部与整体调和分析的含义。
\item 第三卷（第 4 部分）汇集了如今被认为属于相对 Langlands 纲领的主题，从周期与 theta 对应（Howe 对偶）开始，到它们的算术版本结束。Akshay Venkatesh 关于「哈密顿作用与 Langlands 对偶」的讲座——遗憾的是我们未能为本文集收到相应文稿——解释了如何以统一的方式把群 $G$ 上的周期与 theta 级数看作适当哈密顿 $G$-空间的量子化。
\end{itemize}

在准备本文集时，编者并未追求各篇文章在篇幅与范围上的统一，而是让作者自行决定分享多少内容，以他们觉得合适且合时的程度为准。因此，文章篇幅从约 20 页到 200 余页不等。产生这些文章的报告均有录像并公开可得\footnote{见 \url{https://youtube.com/playlist?list=PLx5f8IelFRgEv3qXY43HxVELsFwBKAFhN}。}，但多数文章的内容超出了报告所讲的范围。

暑期学校的一个亮点是 7 月 20 日下午的座谈环节。其间，几位参加过 Corvallis 会议的学者分享了他们对那次会议的回忆，以及 Corvallis 对其后研究生涯的影响。参加者（现场或线上）包括 James Arthur、Bill Casselman、Dick Gross、Michael Harris、Guy Henniart、Hervé Jacquet、Jean-Pierre Labesse、Ken Ribet、Christophe Soulé 与 Marie-France Vignéras。感谢他们充满怀旧与启发的分享，这对在场者是非常大的鼓励。

我们深深感谢所有报告人，感谢他们为完成这项广泛而繁重的任务所付出的时间与心力。我们尤其要提到 Ngô Bảo Châu，他最初提议组织一次 IHES 暑期学校；还要提到 Matt Emerton，他对本项目早期而热情的投入，帮助我们组织起一份星光熠熠的报告人名单，为项目的成功作出了贡献。我们非常感谢 IHES 的工作人员，特别是 Elisabeth Jasserand，感谢他们提供后勤支持，使学校能够接待如此多的参加者，并高效处理一切事务，让我们得以专注于数学。我们感谢赞助方——IHES、NSF 与 Clay 数学研究所——的经费支持。最后，我们感谢 AMS 支持出版，并在各卷的格式与篇幅上给予我们充分的自由。
